\documentclass[11pt]{article}
\usepackage[margin=1in]{geometry}
\usepackage{amsmath,amssymb}
\usepackage[hidelinks]{hyperref}

\newcommand{\Z}{\mathbb Z}
\newcommand{\R}{\mathbb R}

\title{A Literature-Implied Cross-Polytope Subcase of a Small-Radii Maximal-Function Future-Work Direction}
\author{Run fa\_banach\_001}
\date{2026-07-19}

\begin{document}
\maketitle

\paragraph{Status.}
This is a lightweight literature-implied status note:
\[
\texttt{literature\_implied\_answer (partial future-work subcase).}
\]
It records a later answer to the \(q=1\) cross-polytope instance of a
future-work direction in Niksinski--Wrobel,
\emph{Dimension-free estimates for discrete maximal functions and lattice
points in high-dimensional spheres and balls with small radii},
arXiv:2503.16952.

\paragraph{Source direction.}
The source paper defines, in its introduction,
\[
B^q(d)=\{x\in\R^d:\sum_{j=1}^d |x_j|^q\le 1\},\qquad 1\le q<\infty.
\]
In Remark 3 of the introduction, immediately after Theorems 1.1--1.5 and the
discussion of the remaining Stein range, the authors write that their methods
should also extend to maximal functions associated with other \(B^q\) balls,
provided that the supremum is restricted to small radii, and they identify this
as future work.  This packet treats that as a future-work target rather than as
a formal numbered open problem.

\paragraph{Supporting literature.}
Kosz--Niksinski--Wrobel,
\emph{Uniform estimates for Delannoy numbers and dimension-free estimates for
discrete maximal functions over cross-polytopes}, arXiv:2604.15844, studies
the \(q=1\) case.  Its introduction identifies the relevant sets as
cross-polytopes, hyper-octahedrons, or \(\ell^1\) balls, and defines the
discrete averages \(M_R^d\) over
\[
B_R^d=\{x\in\R^d: |x_1|+\cdots+|x_d|\le R\}.
\]
Theorem 1.15 states that, for every \(\varepsilon\in(0,1)\), there is
\(C(\varepsilon)\) such that for all dimensions \(d\) and all \(p\in[2,\infty)\),
\[
\|S^d_{\ast,(0,d^{1-\varepsilon}]} f\|_{\ell^p(\Z^d)}
   \le C(\varepsilon)\|f\|_{\ell^p(\Z^d)}
\]
and consequently
\[
\|M^d_{\ast,(0,d^{1-\varepsilon}]} f\|_{\ell^p(\Z^d)}
   \le C(\varepsilon)\|f\|_{\ell^p(\Z^d)}.
\]
Thus the full maximal function over \(B^1\) balls has a dimension-free
small-radii estimate in the same sense as the future-work direction of
arXiv:2503.16952, at least for the cross-polytope subcase.

The same supporting paper also proves complementary results: Theorem 1.14 gives
dimension-free estimates for large radii \(R\ge c d^{3/2}\) on all
\(\ell^p(\Z^d)\), \(1<p<\infty\), and Theorem 1.17 gives a dimension-free
dyadic maximal estimate for all radii and \(p\in[2,\infty)\).  The paper cites
arXiv:2503.16952 as reference [19] and explicitly says that its proof of
Theorem 1.15 follows the methods of that paper, with the new cross-polytope
lattice-point count as the input.

\paragraph{Scope limitations.}
This is not a new proof produced by this run.  It is a later literature
identification.  It does not answer Stein's original Euclidean-ball question,
does not remove the \(\varepsilon\) loss from arXiv:2503.16952 for Euclidean
balls or spheres, and does not cover all \(q\).  For cross-polytopes, it still
leaves the full non-dyadic middle range between the small radii treated by
Theorem 1.15 and the large radii treated by Theorem 1.14.

\paragraph{Search evidence.}
Local cheap indexes contained no exact prior packet for arXiv:2503.16952 or
arXiv:2604.15844.  A web/arXiv search for arXiv:2604.15844 identified the
cross-polytope paper and its abstract-level claims; the downloaded PDF was then
checked directly for Theorems 1.14, 1.15, 1.17 and for its citation of
arXiv:2503.16952.

\paragraph{References.}
\begin{itemize}
\item J. Niksinski and B. Wrobel, \emph{Dimension-free estimates for discrete
maximal functions and lattice points in high-dimensional spheres and balls with
small radii}, arXiv:2503.16952, v2, 2026.
\item D. Kosz, J. Niksinski, and B. Wrobel, \emph{Uniform estimates for
Delannoy numbers and dimension-free estimates for discrete maximal functions
over cross-polytopes}, arXiv:2604.15844, v1, 2026.
\end{itemize}

\end{document}

